\documentclass[11pt]{article}
\usepackage[margin=1in]{geometry}
\usepackage{amsmath,amssymb,amsthm}
\usepackage[colorlinks=true,urlcolor=blue]{hyperref}
\newtheorem{theorem}{Theorem}
\newtheorem{lemma}[theorem]{Lemma}
\newcommand{\cost}{\operatorname{cost}}
\setlength{\emergencystretch}{4em}
\title{What the dynamic reweighting lemma does and does not supply}
\author{Research note; hand argument}
\date{2 October 2026}
\begin{document}
\maketitle

\noindent\textbf{Scope.}
This note isolates a failed candidate for obtaining the exact
$O(\log\log\log n)$ multi-splay bound from secondary preferred bits.
It does not refute that competitive ratio, a different weight assignment,
or a potential involving the actual splay shapes.

\section{The reweighting theorem uses a lazy accounting state}

Georgakopoulos's reweighting theorem is stated, and its informal proof
reproduced, in Wang's 2006 thesis, Section 1.1.1, pages 3--5.
For a sequence of splays and analytical weight changes, a splay at $x$
can be charged
\[
 c_s\log\frac{W}{w(x)}+O(1),
\]
and increasing $w(x)$ to $w'(x)$ can be charged
\[
 (c_s+1)\log\frac{w'(x)}{w(x)}.
                                                        \tag{1}
\]
Decreases require no positive charge. This is a bound on the total
ledger; it is not the pointwise potential change for eagerly changing
all weights in the ordinary sum-of-log-subtree-masses potential.

For example, give the bottom vertex of an $N$-node chain weight one,
and every other vertex weight $\varepsilon\ll1/N$. Doubling the bottom
weight nearly doubles every ancestor subtree mass. The eager change
of the ordinary potential approaches $N$, although the logarithmic
weight ratio is only one.

The theorem remains valid because its proof keeps a second, lazy weight
assignment. Decreases happen immediately, but an increase is deferred
until the node is next splayed to the root. The eager total mass is an
upper bound on the lazy total mass. Only the first splay of that node
can be undercharged by the eager denominator, and the discrepancy
together with increasing its weight at the root costs at most (1).
Surviving lazy weights and pending increases are accounting state;
they cannot simply be discarded when an auxiliary component is split.

\section{A natural secondary-preference weight assignment}

Let $B$ be a perfect reference on $2^b$ depth leaves
$0,\ldots,H$, where $H=2^b-1$. A primary path has $H$ original
separator keys at depths $0,\ldots,H-1$. For any choice of secondary
preferences, let $a_B(j)$ be the number of nonpreferred edges on the
$B$-path to depth leaf $j$. Fix $\rho=2b$ and assign its separator
at depth $j$ the weight
\[
 w(j)=\rho^{-a_B(j)},\qquad
 \widehat w(j)=\frac{w(j)}{\sum_{i=0}^{H-1}w(i)}.
                                                        \tag{2}
\]
Weights depend on the actual fixed-owner secondary histories of this
path, not on its splay shape. This resembles the interleave-based
weight proposed in Lucier's 2005 project report, applied at the
secondary scale. Its appeal is that a request needing $f$ secondary
changes has a small logarithmic access ratio.

\begin{lemma}[Uniform normalization bounds]\label{lem:norm}
For every assignment of the preferences of a perfect $B$,
\[
 \sum_{j=0}^{H}\rho^{-a_B(j)}=(1+1/\rho)^b.
\]
Consequently the normalizer $W$ in (2) satisfies
\[
 \frac12\le (1+1/\rho)^b-1\le W
       \le(1+1/\rho)^b<2.
                                                        \tag{3}
\]
The ratio of the normalizers of any two path states is less than four.
\end{lemma}

\begin{proof}
At each reference node, one child's generating polynomial is multiplied
by $1$ and the other by $z$. Both subtrees have the same remaining
height. Induction therefore gives $(1+z)^b$, independently of which
child is preferred. Substitute $z=1/\rho$. The routing-key alphabet
omits only depth leaf $H$, whose weight is at most one. Bernoulli's
inequality and $(1+1/(2b))^b\le e^{1/2}<2$ prove (3).
\end{proof}

\begin{theorem}[Recurrent bulk reweighting cost]\label{thm:bulk}
For the normalized weights (2), the sum of positive logarithmic
weight increases on surviving primary-path keys cannot be bounded by
$O((m+S+F)\log b+n\log b)$ over all request sequences.
There is a two-key periodic sequence with $m=2$, $S=3$, and $F=2$
per steady period, whose positive reweighting ledger is
$\Omega(H\log b)$ per period.
\end{theorem}

\begin{proof}
Use a perfect primary node reference of height $H$, initially all right.
Let $r$ be its maximum key, at depth $H-1$. Let $q$ be the node of
the all-right primary spine at depth
\[
 d=2^{b-1}-1.
\]
Alternate the original requests $q,r,q,r,\ldots$. Every proper
ancestor direction other than $q$ stays right. In steady state a
request to $q$ has just its final switch, and a request to $r$ first
switches $q$ back toward the right and then performs its final switch
at the leaf $r$. Thus the actual depth-emission blocks are
\[
 (d),\qquad(d,H-1).                            \tag{4}
\]
There are three switches per period. A leaf's final preference toggle
has no effect on the physical primary-path key partition.

In $B$, leaf $d$ is the rightmost leaf of the left half. Its request
writes left at the root and right at every node below that root.
Leaf $H-1$ is in the right half and writes right at the root. Every
nonroot owner row visited by (4) is always written with the same bit.
After the finite initial writes, the root row therefore contributes
exactly two changes per period and all other rows contribute zero.
In particular $F=2$ in each steady period.

Immediately before a request to $q$, the root bit of $B$ is right.
Immediately afterwards it is left. For each $j\le d$, all its other
secondary path bits are unchanged: those intervals lie in the left
half, and their owner depths lie on the common unchanged primary
prefix through $q$. Hence
\[
 a_B^{\mathrm{after}}(j)=a_B^{\mathrm{before}}(j)-1.
\]
These $d+1=2^{b-1}$ separator keys survive the primary splice; none is
an incoming or deleted tail key. Their raw weight increases by exactly
$\rho$. By Lemma~\ref{lem:norm}, each normalized weight increases by
at least $\rho/4=b/2$. For $b\ge4$, their positive logarithmic
increase is at least
\[
 2^{b-1}\log_2(b/2)=\Omega(H\log b).          \tag{5}
\]
This happens on every request to $q$. The fixed initialization allowance
can be exceeded by repeating the period. Equations (4)--(5) prove the
claimed failure of this particular positive-increase ledger.
\end{proof}

The example is compatible with efficient actual multi-splay behavior:
only two original keys are queried, so its working-set theorem gives
constant amortized cost after initialization. Thus (5) measures wasted
accounting in (2), not a costly multi-splay execution. The obstruction
also survives any treatment that freely rescales an entire component:
the normalized ratios used above already remove that common scale.

The missing ingredient in this route would have to exploit cancellations
between the two halves, actual splay geometry, or a different persistent
weight scheme. Applying the positive-only reweight theorem directly
to the bad-edge weights does not supply those cancellations.
In particular, Theorem~\ref{thm:bulk} does \emph{not} rule out a refined
lazy implementation of the same desired weights. An increase at a cold
prefix key may be cancelled by a later decrease before that key is
splayed. The displayed lower bound concerns the sufficient accounting
rule that charges every eager positive increase. A successful lazy
refinement must separately control the access charges made with its
possibly smaller retained weights, as well as pending state transported
through the primary splits and joins.

\section{A fixed-component limitation of the simplest lazy repair}

The following observation concerns an ordinary fixed component; it is
not an embedding of that component's word into standard multi-splay.
It explains why cancelling cold increases alone does not complete the
argument.

Take a perfect secondary reference $B$ on $N=2^b$ stored keys
$0,\ldots,N-1$, with desired weights $w(x)=\rho^{-a_B(x)}$.
Maintain a lazy weight $\ell(x)$ by applying every desired decrease
immediately, ignoring desired increases until $x$ is accessed, and
then setting $\ell(x)$ to its desired post-access value one. This is
the direct lazy realization suggested by the reweighting proof.

Consider repeated increasing scans $0,1,\ldots,N-1$. Before an access
to $x$, its retained weight is bounded above by the smallest desired
weight it had since its previous access. At some time in that interval
the scan is in the opposite half of $B$. All preferences inside $x$'s
half are then right: that half was completed in the current or the
previous scan. If $z(x)$ counts zero bits in the lower $b-1$ bits of
$x$, it follows that
\[
 \ell(x)\le\rho^{-(1+z(x))}.                  \tag{6}
\]
Immediately before $x$ is accessed, the preceding requested key has
lazy weight one, so the current lazy total mass $W_{\ell}$ is at least
one. The sum of the logarithmic access charges in one steady scan is
therefore at least
\[
 \sum_{x=0}^{N-1}\log\frac{W_{\ell}}{\ell(x)}
 \ge \log\rho\sum_{x=0}^{N-1}(1+z(x))
 =\frac{N(b+1)}2\log\rho.                    \tag{7}
\]
On the other hand each secondary reference vertex changes preference
exactly twice per cyclic scan, so $F=2(N-1)$. For $\rho=2b$, the
charge in (7) is $\Omega(Nb\log b)$ although
$(N+F)\log b=O(N\log b)$. Ordinary splay itself has linear scan cost.

Thus the direct per-node lazy minimum can badly overcharge accesses
after it has saved the eager reweight charges. A multi-splay-specific
repair would have to exploit its component cuts, rescaling, or additional
state; the fixed-component calculation does not rule out such a repair.

\paragraph{Sources.}
The reweighting statement and lazy proof are in
\href{https://www.csd.cmu.edu/sites/default/files/phd-thesis/CMU-CS-06-140.pdf}
{Chengwen Chris Wang, \emph{Multi-Splay Trees} (2006)},
Section 1.1.1; the working-set result follows from its remassing lemma
in Section 3.4.3. The original reweighting source is
\href{https://doi.org/10.1016/j.jalgor.2003.09.004}
{G. F. Georgakopoulos, \emph{Splay trees: a reweighing lemma and a proof
of competitiveness vs. dynamic balanced trees} (2004)}.
The related secondary-depth weight idea is motivated by Brendan Lucier,
\emph{Applying the Interleave Bound to Splay Trees}, project report,
version 0.1, 13 April 2005; the present obstruction concerns (2), not
all conclusions or conjectures of that report.

\end{document}
